\chapter{Finsler manifolds as metric spaces}
\label{chap:metric-spaces}

A Finsler structure turns a manifold into a metric space, and the question is how much of the
manifold can be recovered from the distance. This chapter follows Part~II of \cite{Lar22}:
Chapter~8 (interior metric spaces, after Gromov) and Chapter~9 (Finsler manifolds). The first
section works in an arbitrary metric space; the second identifies $(M, d_b)$ as an interior
metric space; the third compares the topology of $d_b$ with the manifold topology, which is where
the model space enters.

\section{The rectifiable distance}
\label{sec:ms-rectifiable}

Let $(X, d)$ be a pseudo-metric space, with distances in $[0, \infty]$. The length of a continuous
curve $\gamma \colon [a, b] \to X$ is the supremum of the sums
$\sum_i d(\gamma(t_i), \gamma(t_{i+1}))$ over the partitions of $[a, b]$; in Mathlib this is the
variation \texttt{eVariationOn}. Its basic properties are in Mathlib: refining a partition
increases the sums, the length vanishes only for constant curves, it is additive under
concatenation, invariant under monotone reparametrizations \cite[Obs.~8.1.1]{Lar22}, lower
semicontinuous under uniform convergence \cite[Prop.~8.1.3]{Lar22}, and the arc length
$t \mapsto \ell(\gamma|_{[a, t]})$ is continuous \cite[Lema~8.1.4]{Lar22}.

\begin{definition}[The rectifiable distance]
  \label{def:ms-length-dist}
  \lean{LCManifolds.RectPath, LCManifolds.RectPath.length, LCManifolds.RectPath.concat, LCManifolds.RectPath.reverse, LCManifolds.RectPath.length_concat, LCManifolds.RectPath.length_reverse, LCManifolds.lengthDist}
  \leanok
  The rectifiable distance $d_\ell(x, y)$ is the infimum of the lengths of the continuous curves
  from $x$ to $y$, and $\infty$ if there is none.
\end{definition}

\begin{proposition}
  \label{prop:ms-length-dist}
  \lean{LCManifolds.lengthDist_self, LCManifolds.lengthDist_comm, LCManifolds.lengthDist_triangle, LCManifolds.edist_le_lengthDist, LCManifolds.LengthMetric, LCManifolds.lipschitzWith_ofLength, LCManifolds.continuous_ofLength}
  \leanok
  \uses{def:ms-length-dist}
  $d_\ell$ is a pseudo-distance, and $d \leq d_\ell$. In particular the topology of $d_\ell$ is
  finer than that of $d$ \cite[Obs.~8.1.7]{Lar22}.
\end{proposition}

\begin{proof}
  \leanok
  Constant curves, reversed curves and concatenations give the three properties of a distance.
  The trivial partition gives $d(x, y) \leq \ell(\gamma)$.
\end{proof}

\begin{lemma}[When the topologies agree]
  \label{lem:ms-topologies}
  \lean{LCManifolds.continuous_toLength_iff}
  \leanok
  \uses{prop:ms-length-dist}
  The topologies of $d$ and $d_\ell$ coincide if and only if, for every $x$ and every
  $\varepsilon > 0$, the points close enough to $x$ can be joined to $x$ by curves of length
  $< \varepsilon$ \cite[Lema~8.1.8]{Lar22}.
\end{lemma}

\begin{proof}
  \leanok
  This is the continuity of the identity $(X, d) \to (X, d_\ell)$, written with balls.
\end{proof}

\begin{definition}[Interior metric spaces]
  \label{def:ms-interior}
  \lean{LCManifolds.IsInterior, LCManifolds.IsInterior.exists_rectPath}
  \leanok
  \uses{def:ms-length-dist}
  $(X, d)$ is an interior metric space (a length space) if $d = d_\ell$
  \cite[\S8.1.1]{Lar22}.
\end{definition}

\begin{proposition}[Approximate midpoints]
  \label{prop:ms-midpoints}
  \lean{LCManifolds.IsInterior.exists_approx_midpoint}
  \leanok
  \uses{def:ms-interior}
  In an interior metric space, for all $x, y$ and $\varepsilon > 0$ there is $z$ with
  $\max\{d(x, z), d(z, y)\} \leq \frac12 d(x, y) + \varepsilon$ \cite[Problema~8.III]{Lar22}.
\end{proposition}

\begin{proof}
  \leanok
  Take a curve of length $< d(x, y) + 2\varepsilon$. By the continuity of the arc length and the
  intermediate value theorem it has a point at half its length.
\end{proof}

\section{The metric Hopf--Rinow theorem}
\label{sec:ms-hopf-rinow}

\begin{lemma}[Approximate intermediate points]
  \label{lem:ms-intermediate}
  \lean{LCManifolds.IsInterior.exists_dist_le_add, LCManifolds.IsInterior.exists_approx_intermediate, LCManifolds.IsInterior.exists_mem_closedBall_near}
  \leanok
  \uses{prop:ms-midpoints}
  In an interior metric space, for $a, b \geq 0$ with $a + b = d(x, y)$ and $\varepsilon > 0$
  there is $z$ with $d(x, z) \leq a + \varepsilon$ and $d(z, y) \leq b + \varepsilon$.
\end{lemma}

\begin{proof}
  \leanok
  As in Proposition~\ref{prop:ms-midpoints}, with the point at arc length $a$.
\end{proof}

\begin{lemma}[Lipschitz reparametrization]
  \label{lem:ms-reparam}
  \lean{LCManifolds.RectPath.exists_lipschitz_reparam}
  \leanok
  \uses{def:ms-length-dist}
  A curve of finite length $\ell$ can be reparametrized on $[0, 1]$ as an $\ell$-Lipschitz curve
  \cite[Obs.~8.1.5, Cor.~8.1.6]{Lar22}.
\end{lemma}

\begin{proof}
  \leanok
  Compose with a right inverse of the normalized arc length; the curve is constant on the level
  sets of the arc length.
\end{proof}

\begin{theorem}[\texorpdfstring{Hopf--Rinow \cite[Teorema~8.2.3]{Lar22}}{Hopf-Rinow}]
  \label{thm:ms-hopf-rinow}
  \lean{LCManifolds.IsInterior.isCompact_closedBall_of_lt, LCManifolds.IsInterior.exists_isCompact_closedBall_gt, LCManifolds.IsInterior.isCompact_closedBall, LCManifolds.IsInterior.properSpace, LCManifolds.IsInterior.exists_rectPath_length_eq_of_isCompact, LCManifolds.IsInterior.exists_rectPath_length_eq}
  \leanok
  \uses{lem:ms-intermediate, lem:ms-reparam}
  Let $X$ be a complete interior metric space.
  \begin{enumerate}
    \item If $\bar B_r(x)$ is compact for all $0 < r < \rho$, then $\bar B_\rho(x)$ is compact.
    \item If $X$ is locally compact, every closed ball is compact, and any two points are joined
      by a shortest curve, of length equal to their distance.
  \end{enumerate}
\end{theorem}

\begin{proof}
  \leanok
  (1) By Lemma~\ref{lem:ms-intermediate} every point of $\bar B_\rho(x)$ is close to the compact
  ball $\bar B_{\rho - \delta}(x)$, so $\bar B_\rho(x)$ is totally bounded, and it is complete.
  (2) The supremum of the radii of compact balls cannot be finite: by (1) the ball of that radius
  is compact, and a finite cover by balls with compact closure, together with a Lebesgue number,
  gives a larger compact ball. For the shortest curves, reparametrize almost minimizing curves as
  in Lemma~\ref{lem:ms-reparam}; they lie in a compact ball and are uniformly Lipschitz, and a
  cluster point in the product topology (Tychonoff) is a Lipschitz curve of length at most the
  distance. Here the formal proof departs from \cite{Lar22}, which uses Arzelà--Ascoli and the
  lower semicontinuity of the length.
\end{proof}

\begin{remark}
  \label{rem:ms-proper}
  \lean{LCManifolds.completeSpace_of_isCompact_closedBall}
  \leanok
  A space whose closed balls are compact is complete \cite[Obs.~8.2.4]{Lar22}. Being geodesic and
  locally compact does not imply completeness: the interval $(0, 1)$.
\end{remark}

\section{\texorpdfstring{$(M, d_b)$}{(M, d)} is an interior metric space}
\label{sec:ms-finsler-interior}

Let $\mu$ be a reversible Finsler length structure on $M$ (Definition~\ref{def:finsler-length-structure}).
Its distance $d_b$ (Definition~\ref{def:mt-dist}) is then a pseudo-distance, and $(M, d_b)$ is a
pseudo-metric space; without reversibility $d_b$ is only a quasi-distance and the symmetric
theory of the previous section does not apply as such.

\begin{theorem}[\texorpdfstring{\cite[Prop.~9.1.2]{Lar22}}{Prop. 9.1.2}]
  \label{thm:ms-finsler-interior}
  \lean{LCManifolds.FinslerLengthStructure.FinslerSpace, LCManifolds.FinslerLengthStructure.edist_le_abs_length, LCManifolds.FinslerLengthStructure.continuous_toFinsler, LCManifolds.FinslerLengthStructure.eVariationOn_le_length, LCManifolds.FinslerLengthStructure.arcPath, LCManifolds.FinslerLengthStructure.lengthDist_le_of_reach, LCManifolds.FinslerLengthStructure.isInterior}
  \leanok
  \uses{prop:ms-length-dist, def:ms-interior, thm:mt-triangle, prop:mt-dist-comm}
  $\Ck{1}$ curves are continuous for $d_b$, the metric length of a $\Ck{1}$ arc is at most its
  Finsler length, $\ell(\gamma) \leq L_b(\gamma)$, and $(M, d_b)$ is an interior metric space.
\end{theorem}

\begin{proof}
  \leanok
  $d_b(\gamma(s), \gamma(t)) \leq |L_b(\gamma|_{[s, t]})|$ gives the continuity. For a partition,
  $\sum d_b(\gamma(t_i), \gamma(t_{i+1})) \leq \sum L_b(\gamma|_{[t_i, t_{i+1}]}) \leq L_b(\gamma)$.
  Hence $d_\ell$ is bounded by the length of every chain of $\Ck{1}$ arcs, so $d_\ell \leq d_b$;
  the other inequality is Proposition~\ref{prop:ms-length-dist}.
\end{proof}

\section{Bounded Finsler structures and the topology}
\label{sec:ms-bounded}

In this section the model $E$ is a normed space; from Lemma~\ref{lem:ms-exit} on, it is complete
and $M$ is Hausdorff.

\begin{definition}[\texorpdfstring{\cite[Def.~9.1.3, Obs.~9.1.4]{Lar22}}{Def. 9.1.3}]
  \label{def:ms-bounded}
  \lean{LCManifolds.FinslerSeminormField.IsBoundedBelowOn, LCManifolds.FinslerSeminormField.IsBoundedAboveOn, LCManifolds.FinslerSeminormField.IsBoundedBelow, LCManifolds.FinslerSeminormField.IsBoundedAbove, LCManifolds.FinslerSeminormField.IsBounded, LCManifolds.FinslerSeminormField.IsBounded.isBoundedBelow, LCManifolds.FinslerSeminormField.IsBounded.isBoundedAbove}
  \leanok
  $\mu$ is \emph{bounded below} if every point has an open neighborhood $U$ in the domain of a
  chart $\varphi$ and $C > 0$ with $\|\varphi_{*x} v\| \leq C \mu_x(v)$ for $x \in U$; it is
  \emph{bounded above} if $m \mu_x(v) \leq \|\varphi_{*x} v\|$ for some $m > 0$. A Finsler
  structure bounded above gives a \emph{weak Finsler manifold}; bounded above and below, a
  \emph{strong} one.
\end{definition}

\begin{remark}[Locally convex models]
  If the model is not normable, no Finsler norm $\mu_x$ can be bounded above and below by a norm
  inducing the topology of the model, since that topology would then be a norm topology. Strong
  Finsler manifolds are a phenomenon of normed models.
\end{remark}

\begin{lemma}[The integrated bound]
  \label{lem:ms-integrated}
  \lean{LCManifolds.FinslerLengthStructure.norm_chart_sub_le}
  \leanok
  \uses{def:ms-bounded}
  If $\mu$ is bounded below on $U$ with constant $C$ and $E$ is complete, a $\Ck{1}$ curve that
  stays in $U$ satisfies $\|\varphi(\gamma(t)) - \varphi(\gamma(s))\| \leq C L_b(\gamma|_{[s, t]})$
  \cite[(9.2)]{Lar22}.
\end{lemma}

\begin{proof}
  \leanok
  Integrate the derivative of $\varphi \circ \gamma$ and bound its norm by $C \mu(\dot\gamma)$.
\end{proof}

\begin{lemma}[Leaving a chart costs length]
  \label{lem:ms-exit}
  \lean{LCManifolds.FinslerLengthStructure.arc_exit, LCManifolds.FinslerLengthStructure.reach_exit}
  \leanok
  \uses{lem:ms-integrated, def:mt-chain}
  Let $\mu$ be bounded below on $U$ with constant $C$, and let $\bar B(\varphi(x_0), r)$ be the
  image of a part of $U$. A chain of $\Ck{1}$ arcs of length $L$ starting at $x_0$ either has
  $C L \geq r$, or ends at a point $y \in U$ with $\|\varphi(y) - \varphi(x_0)\| \leq C L$.
\end{lemma}

\begin{proof}
  \leanok
  For a single arc, let $t_0$ be the first time at which the arc leaves
  $U \cap \varphi^{-1}(B(\varphi(x_0), r))$. If $C L_b(\gamma|_{[a, t_0]}) < r$, the points
  $\varphi(\gamma(t))$, $t < t_0$, form a Cauchy family in a closed ball inside $\varphi(U)$; by
  completeness they converge, and since $M$ is Hausdorff the limit is $\varphi(\gamma(t_0))$, which
  is then inside the ball: a contradiction. Chains follow by induction.
\end{proof}

\begin{theorem}[\texorpdfstring{\cite[Lema~9.1.8]{Lar22}}{Lema 9.1.8}]
  \label{thm:ms-bounded-below}
  \lean{LCManifolds.FinslerLengthStructure.exists_chart_norm_le_dist, LCManifolds.FinslerLengthStructure.eq_of_dist_eq_zero, LCManifolds.FinslerLengthStructure.continuous_ofFinsler}
  \leanok
  \uses{lem:ms-exit, thm:ms-finsler-interior}
  If $\mu$ is bounded below, every point $x_0$ has a chart and $C, r > 0$ such that
  $d_b(x_0, y) < r / C$ implies that $y$ is in the chart and
  $\|\varphi(y) - \varphi(x_0)\| \leq C d_b(x_0, y)$. Consequently $d_b$ separates points, and, for
  a reversible $\mu$, every open set of $M$ is open for $d_b$.
\end{theorem}

\begin{proof}
  \leanok
  Apply Lemma~\ref{lem:ms-exit} to chains whose length approaches $d_b(x_0, y)$.
\end{proof}

\begin{theorem}[Weak and strong Finsler manifolds]
  \label{thm:ms-bounded-above}
  \lean{LCManifolds.FinslerLengthStructure.dist_le_norm_chart, LCManifolds.FinslerLengthStructure.continuous_toFinsler_of_isBoundedAbove, LCManifolds.FinslerLengthStructure.homeomorphFinsler}
  \leanok
  \uses{def:ms-bounded, thm:ms-bounded-below}
  If $\mu$ is bounded above with constant $m$, then $d_b(x, y) \leq \|\varphi(y) - \varphi(x)\| / m$
  near each point \cite[(9.3)]{Lar22}, so the topology of $d_b$ is coarser than that of $M$. For a
  strong Finsler manifold (reversible, complete model, $M$ Hausdorff) the identity is a
  homeomorphism between $M$ and $(M, d_b)$ \cite[Cor.~9.1.9]{Lar22}.
\end{theorem}

\begin{proof}
  \leanok
  The segment from $\varphi(x)$ to $\varphi(y)$, reparametrized by a smooth transition function so
  that it is a $\Ck{1}$ curve on all of $\R$, has Finsler length at most
  $\|\varphi(y) - \varphi(x)\| / m$. The second statement combines this with
  Theorem~\ref{thm:ms-bounded-below}.
\end{proof}

\begin{lemma}[\texorpdfstring{\cite[Lema~9.1.6]{Lar22}}{Lema 9.1.6}]
  \label{lem:ms-lipschitz}
  \lean{LCManifolds.FinslerSeminormField.chartFun, LCManifolds.FinslerSeminormField.chartFun_tcc, LCManifolds.FinslerSeminormField.chartFun_smul, LCManifolds.FinslerSeminormField.chartFun_zero, LCManifolds.FinslerSeminormField.isBoundedAbove_of_lipschitz}
  \leanok
  \uses{def:ms-bounded}
  If $\mu$, read in charts, is Lipschitz near the zero section, then $\mu$ is bounded above.
\end{lemma}

\begin{proof}
  \leanok
  Near $(\varphi(p), 0)$ the Lipschitz condition gives
  $\mu_q(w) \leq K(\|\varphi(q) - \varphi(p)\| + \|\varphi_{*q} w\|)$; applied to a multiple of $w$ of
  fixed small norm and combined with homogeneity, it gives $\mu_q(w) \leq 2K \|\varphi_{*q} w\|$
  for $q$ close to $p$.
\end{proof}

\begin{theorem}[\texorpdfstring{\cite[Teorema~9.1.5]{Lar22}}{Teorema 9.1.5}]
  \label{thm:ms-finite-dim}
  \lean{LCManifolds.tangentCoordChange_smul, LCManifolds.FinslerMetric.isBounded_of_finiteDimensional}
  \leanok
  \uses{def:ms-bounded, def:finsler-metric}
  If the model is finite-dimensional, every Finsler metric whose seminorms are norms is strong.
\end{theorem}

\begin{proof}
  \leanok
  In a chart, $(y, v) \mapsto \mu(y, v)$ is continuous on the tangent bundle, hence attains a
  positive minimum and a maximum on the compact set $\bar B \times S$, closed ball times unit
  sphere. Homogeneity gives the bounds for every vector.
\end{proof}

\begin{corollary}[\texorpdfstring{Hopf--Rinow for Finsler manifolds \cite[Cor.~9.1.10]{Lar22}}{Hopf-Rinow for Finsler manifolds}]
  \label{cor:ms-finsler-hopf-rinow}
  \lean{LCManifolds.FinslerMetric.metricSpace, LCManifolds.FinslerMetric.locallyCompactSpace, LCManifolds.FinslerMetric.exists_rectPath_length_eq}
  \leanok
  \uses{thm:ms-finite-dim, thm:ms-bounded-above, thm:ms-finsler-interior, thm:ms-hopf-rinow}
  Let $M$ be a Hausdorff manifold modeled on a finite-dimensional normed space, with a reversible
  Finsler metric whose seminorms are norms and with finite distance $d_b$. If $(M, d_b)$ is
  complete, any two points are joined by a continuous curve whose length is their distance.
\end{corollary}

\begin{proof}
  \leanok
  By Theorem~\ref{thm:ms-finite-dim} the metric is strong; by Theorems~\ref{thm:ms-bounded-below}
  and~\ref{thm:ms-bounded-above}, $d_b$ is a distance and $(M, d_b)$ is homeomorphic to $M$, hence
  locally compact. It is an interior space by Theorem~\ref{thm:ms-finsler-interior}, and
  Theorem~\ref{thm:ms-hopf-rinow} applies.
\end{proof}

Whether these shortest curves can be taken piecewise $\Ck{1}$ is a separate question
\cite[Obs.~9.1.11]{Lar22}.

\section{Asymmetric distances}
\label{sec:ms-asymmetric}

For a non-reversible Finsler structure $d_b$ is not symmetric. Following Mennucci \cite{Men13},
an asymmetric distance $b$ satisfies $b(x, x) = 0$ and the triangle inequality, and the topology is
that of the symmetrized distance $\max\{b(x, y), b(y, x)\}$; forward and backward balls alone may
generate different topologies \cite[Example~4.1]{Men13}. Lengths are forward lengths: suprema of
$\sum b(\gamma(t_i), \gamma(t_{i+1}))$, which change under reversal of the curve
\cite[Remark~3.5]{Men13}.

\begin{definition}[Asymmetric distances and forward lengths]
  \label{def:ms-asymm}
  \lean{LCManifolds.AsymmPseudoEMetricSpace, LCManifolds.AsymmPseudoEMetricSpace.ofBDist, LCManifolds.fwdVar, LCManifolds.bwdVar, LCManifolds.FwdPath, LCManifolds.fwdLengthDist, LCManifolds.IsFwdInterior}
  \leanok
  An asymmetric metric space carries $b$ and the topology of its symmetrization. The forward
  length of a curve is its forward variation; $b_\ell$ is the infimum of forward lengths of
  continuous curves, and $X$ is forward interior if $b_\ell = b$ (g-intrinsic in
  \cite[Def.~3.16]{Men13}).
\end{definition}

\begin{theorem}[Forward lengths]
  \label{thm:ms-fwd-length}
  \lean{LCManifolds.fwdVar.Icc_add_Icc, LCManifolds.fwdVar.comp_eq_of_monotoneOn, LCManifolds.fwdVar_comp_neg, LCManifolds.fwdVar.continuousOn_Icc, LCManifolds.fwdLengthDist_triangle, LCManifolds.bdist_le_fwdLengthDist, LCManifolds.IsFwdInterior.exists_bdist_le_add, LCManifolds.FwdPath.exists_lipschitz_reparam}
  \leanok
  \uses{def:ms-asymm}
  Forward lengths are additive and invariant under increasing reparametrizations; the running
  length of a continuous curve of finite forward length is continuous \cite[Prop.~3.9]{Men13};
  $b \leq b_\ell$ and $b_\ell$ satisfies the triangle inequality. Forward interior spaces have
  approximate intermediate points, and curves of finite forward length admit forward Lipschitz
  reparametrizations on $[0, 1]$.
\end{theorem}

\begin{proof}
  \leanok
  The proofs of the symmetric case adapt, with care for the order of the arguments of $b$; the
  continuity of the running length follows \cite[Appendix~B.2]{Men13}.
\end{proof}

\begin{theorem}[Shortest forward curves]
  \label{thm:ms-fwd-hopf-rinow}
  \lean{LCManifolds.bdist_le_bdist_add_edist, LCManifolds.continuous_bdist, LCManifolds.IsFwdInterior.exists_fwdPath_length_eq_of_isCompact}
  \leanok
  \uses{thm:ms-fwd-length}
  Let $X$ be a Hausdorff forward interior space, $b(x, y) < \infty$, with compact closed forward
  ball $\{z : b(x, z) \leq b(x, y) + 1\}$ on which $b(w, z) \leq C\, b(z, w)$. Then there is a
  continuous curve from $x$ to $y$ of forward length $b(x, y)$.
\end{theorem}

\begin{proof}
  \leanok
  As in Theorem~\ref{thm:ms-hopf-rinow}, with forward Lipschitz reparametrizations. The bounded
  asymmetry makes the limit curve Lipschitz for the symmetrized distance, hence continuous; without
  it the limit need not be continuous \cite{Men13}.
\end{proof}

\begin{theorem}[Non-reversible Finsler structures]
  \label{thm:ms-finsler-asymm}
  \lean{LCManifolds.FinslerLengthStructure.bdist_le_lengths, LCManifolds.FinslerLengthStructure.fwdVar_le_length, LCManifolds.FinslerLengthStructure.fwdArcPath, LCManifolds.FinslerLengthStructure.fwdLengthDist_le_of_reach, LCManifolds.FinslerLengthStructure.isFwdInterior, LCManifolds.FinslerLengthStructure.length_reverse_le, LCManifolds.FinslerLengthStructure.Reach.reverse_le, LCManifolds.FinslerLengthStructure.dist_comm_le, LCManifolds.FinslerLengthStructure.t2Space_finslerSpace, LCManifolds.FinslerLengthStructure.exists_fwdPath_length_eq}
  \leanok
  \uses{thm:ms-fwd-hopf-rinow, thm:ms-finsler-interior, thm:ms-bounded-below}
  For any Finsler length structure, $\ell^+(\gamma) \leq L_b(\gamma)$ and $(M, d_b)$ is forward
  interior. A reversibility constant $\mu_x(-v) \leq \lambda \mu_x(v)$ gives
  $d_b(y, x) \leq \lambda d_b(x, y)$; if moreover $\mu$ is bounded below and the closed forward
  ball is compact, there are shortest forward curves.
\end{theorem}

\begin{proof}
  \leanok
  As Theorem~\ref{thm:ms-finsler-interior}, with forward sums; reversed chains give the bound on
  the asymmetry, and Lemma~\ref{thm:ms-bounded-below} makes $(M, d_b)$ Hausdorff.
\end{proof}
